% Part: proof-theory
% Chapter: cut-elimination
% Section: interpolation

\documentclass[../../../include/open-logic-section]{subfiles}

\begin{document}

\olfileid[pt]{pt}{cut}{itp}

\olsection{Lema de Maehara e teorema da interpolação de Craig}

O teorema da interpolação, devido a William Craig, afirma que, se $!A
\Entails !B$, então existe !!a{sentence}~$!C$ (o \emph{interpolante}) tal
que $!A \Entails !C$ e $!C \Entails !B$. É importante que o
interpolante~$!C$ possa ser escolhido de modo que todas as !!{constant}s
e todos os !!{predicate}s que ocorrem em~$!C$ ocorram tanto em~$!A$
quanto em~$!B$. (Supomos que não haja !!{function}s envolvidas.) Sem
essa restrição, $!A$ e $!B$ poderiam servir trivialmente como
interpolantes.

O teorema da interpolação vale para a lógica clássica de primeira ordem
e já foi chamado de ``a última propriedade importante da lógica de
primeira ordem'' a ser descoberta. Ele tem aplicações importantes na
teoria de modelos e no raciocínio automatizado. Sua motivação e aplicação
originais estavam na filosofia da ciência, pois ele pode ser usado para
estudar os conceitos primitivos mediante os quais uma teoria pode ou não
ser axiomatizada. Ele também implica o teorema da definibilidade de Beth,
segundo o qual, se uma teoria pode definir implicitamente um conceito,
também pode defini-lo explicitamente.

Como muitos resultados da lógica matemática, o teorema da interpolação
pode ser demonstrado tanto pelos métodos da teoria de modelos quanto
pelos da teoria da prova. A vantagem dos métodos da teoria da prova é
que eles não apenas mostram a existência de um interpolante, mas também
fornecem um algoritmo que o calcula a partir de !!a{proof} de que $!A
\Entails !B$, por exemplo, !!a{proof} em~\Log{G3c} de $!A \Sequent !B$.

Seja $\Lang L(!A)$ o conjunto das !!{constant}s e dos !!{predicate}s
em~$!A$, e seja $\Lang L(\Gamma) = \bigcup \Setabs{\Lang L(!A)}{!A \in
\Gamma}$. Ao longo desta seção, supomos que estamos tratando de
!!{formula}s sem !!{function}s.

\begin{thm}[Teorema da interpolação de Craig]\ollabel{thm:interpolation}
Seja a !!{language} de~$!A$ igual a $\Lang L(!A)$, e a de $!B$, igual a
$\Lang L(!B)$.
Se $\Log{G3c} \Proves !A \Sequent !B$, então existe !!a{formula}~$!C$ na
!!{language}~$\Lang L(!C) \subseteq \Lang L(!A) \cap \Lang L(!B)$ tal
que $\Log{G3c} \Proves !A \Sequent !C$ e $\Log{G3c} \Proves !C
\Sequent !B$.
\end{thm}

Demonstramos uma generalização do teorema da interpolação para aproveitar
a estrutura das !!{proof}s em~\Log{G3c}. Para esse fim, imagine que o
antecedente e o sucedente de cada sequente estejam separados em duas
partes. Escrevemos $\maeh{\Gamma_1}{\Gamma_2} \Sequent
\maeh{\Delta_1}{\Delta_2}$ para indicar que $\Gamma_1$ e $\Delta_1$
pertencem à primeira parte, e $\Gamma_2$ e $\Delta_2$, à segunda. O
teorema da interpolação decorre então imediatamente do lema seguinte.

\begin{lem}[Lema de Maehara]\ollabel{lem:maehara} Se $\Log{G3c}
\Proves \maeh{\Gamma_1}{\Gamma_2} \Sequent \maeh{\Delta_1}{\Delta_2}$,
então existe !!a{formula}~$!C$ na !!{language}~$\Lang L(!C)
\subseteq \Lang L(\Gamma_1,
\Delta_1) \cap \Lang L(\Gamma_2, \Delta_2)$ tal que $\Log{G3c}
\Proves \Gamma_1 \Sequent \Delta_1, !C$ e $\Log{G3c} \Proves !C,
\Gamma_2 \Sequent \Delta_2$.
\end{lem}

\begin{proof}
  Suponha que $\pi \Proves \maeh{\Gamma_1}{\Gamma_2} \Sequent
  \maeh{\Delta_1}{\Delta_2}$. Demonstramos a afirmação por indução em
  $n=\pheight{\pi}$.

  Se $n=0$, então $\maeh{\Gamma_1}{\Gamma_2} \Sequent
  \maeh{\Delta_1}{\Delta_2}$ é um axioma, e alguma !!{formula}
  atômica~$!A$ ocorre tanto em $\Gamma_1, \Gamma_2$ quanto em $\Delta_1,
  \Delta_2$. Há quatro casos, conforme $!A$ pertença à primeira ou à
  segunda parte à esquerda, e à primeira ou à segunda parte à direita.
  \begin{enumerate}
    \item $!A \in \Gamma_1$ e $!A \in \Delta_1$. Queremos um~$!C$ tal
    que tanto $\Gamma_1 \Sequent \Delta_1, !C$ quanto $!C, \Gamma_2
    \Sequent \Delta_2$ sejam !!{provable}. O primeiro já é um axioma,
    qualquer que seja a escolha de~$!C$, pois $!A$ ocorre à esquerda e à
    direita. Nossa única opção para garantir que $!C, \Gamma_2 \Sequent
    \Delta_2$ seja !!{provable} é tomar $!C \ident \lfalse$.
    \item $!A \in \Gamma_1$ e $!A \in \Delta_2$. Então $\Gamma_1
    \Sequent \Delta_1, !A$ e $!A, \Gamma_2 \Sequent \Delta_2$ são ambos
    axiomas; portanto, podemos escolher $!C \ident !A$.
    \item $!A \in \Gamma_2$ e $!A \in \Delta_1$. Então $!A, \Gamma_1
    \Sequent \Delta_1$ e $\Gamma_2 \Sequent \Delta_2, !A$ seriam
    axiomas, mas têm~$!A$ nos lados errados. Portanto, $!A$ não é o
    interpolante. Contudo, usando $\RightR{\lnot}$ e $\LeftR{\lnot}$,
    podemos !!{prove} $\Gamma_1 \Sequent \Delta_1, \lnot!A$ e $\lnot
    !A, \Gamma_2 \Sequent \Delta_2$.
    \item $!A \in \Gamma_2$ e $!A \in \Delta_2$. Exercício.
  \end{enumerate}
  Os casos em que $\maeh{\Gamma_1}{\Gamma_2} \Sequent
  \maeh{\Delta_1}{\Delta_2}$ é um axioma porque $\lfalse \in \Gamma_1$
  ou $\lfalse \in \Gamma_2$ ficam como exercícios.

  Se $n>0$, $\pi$ termina em uma inferência lógica. Devemos distinguir
  os casos conforme a regra usada e a parte à qual pertence a fórmula
  principal. Em cada caso, podemos usar a hipótese de indução de que o
  lema vale para todas as !!{proof}s~$\pi_1$ com $\pheight{\pi_1} < n$,
  qualquer que seja a partição em duas partes do sequente final de
  $\pi_1$. Consideremos alguns casos específicos.
  \begin{enumerate}
    \item $\pi$ termina em $\RightR{\lnot}$ com fórmula principal~$\lnot
    !A$. Há dois casos: conforme $\lnot !A$ pertença à primeira ou à
    segunda parte, $\pi$ termina em
    \[
    \AxiomC{}
    \RightLabel{$\pi_1$}
    \Deduce$\maeh{!A, \Gamma_1}{\Gamma_2} \fCenter
      \maeh{\Delta_1}{\Delta_2}$
    \RightLabel{\RightR{\lnot}}
    \UnaryInf$\maeh{\Gamma_1}{\Gamma_2} \fCenter
      \maeh{\Delta_1, \lnot !A}{\Delta_2}$
    \DisplayProof
    \]
    ou
    \[
    \AxiomC{}
    \RightLabel{$\pi_1$}
    \Deduce$\maeh{\Gamma_1}{!A, \Gamma_2} \fCenter
      \maeh{\Delta_1}{\Delta_2}$
    \RightLabel{\RightR{\lnot}}
    \UnaryInf$\maeh{\Gamma_1}{\Gamma_2} \fCenter
      \maeh{\Delta_1}{\Delta_2, \lnot !A}$
    \DisplayProof
    \]
    Observe que, se a !!{formula} principal~$\lnot !A$ pertence à
    primeira parte, atribuímos a !!{formula} ativa~$!A$ também à
    primeira parte da premissa. Essa escolha está a nosso alcance. A
    hipótese de indução também se aplicaria se atribuíssemos a fórmula
    ativa à outra parte, mas então não conseguiríamos encontrar um
    interpolante (exercício: tente fazê-lo).

    Considere primeiro a situação em que $\lnot !A$ pertence à primeira
    parte. Pela hipótese de indução, há um interpolante do sequente final
    de~$\pi_1$, isto é, !!a{formula}~$!C_1$ tal que ambos
    \begin{align*}
    !A, \Gamma_1 & \Sequent \Delta_1, !C_1 &&\text{e} &
    !C_1, \Gamma_2 & \Sequent \Delta_2
    \intertext{são demonstráveis. Precisamos de !!a{formula}~$!C$ tal que
    os dois sequentes seguintes sejam !!{provable}:}
    \Gamma_1 & \Sequent \Delta_1, \lnot !A, !C &&\text{e} &
    !C, \Gamma_2 & \Sequent \Delta_2
    \end{align*}
    É claro que podemos simplesmente tomar $!C \ident !C_1$: a hipótese
    de indução já nos disse que o sequente à direita é !!{provable}, e o
    sequente à esquerda decorre, por~$\RightR{\lnot}$, do sequente obtido
    pela hipótese de indução. A hipótese de indução também nos diz que
    $\Lang L(!C_1) \subseteq \Lang L(!A, \Gamma_1, \Delta_1) \cap \Lang
    L(\Gamma_2, \Delta_2)$. Como $\Lang L(\lnot !A) = \Lang L(!A)$ e
    $\Lang L(!C) = \Lang L(!C_1)$, temos $\Lang L(!C_1) \subseteq \Lang
    L(\Gamma_1, \Delta_1, \lnot !A) \cap \Lang L(\Gamma_2, \Delta_2)$.

    No segundo cenário, a situação é um pouco diferente, pois $\lnot
    !A$ pertence à segunda parte. Também atribuímos à segunda parte a
    !!{formula} ativa da premissa e aplicamos a hipótese de indução a
    \begin{align*}
    \Gamma_1 & \Sequent \Delta_1, !C_1 &&\text{e} &
    !C_1, !A, \Gamma_2 & \Sequent \Delta_2
    \intertext{que são demonstráveis. Se aplicarmos novamente a regra \RightR{\lnot} ao segundo deles, obteremos !!{proof}s de}
    \Gamma_1 & \Sequent \Delta_1, !C_1 &&\text{e} &
    !C_1, \Gamma_2 & \Sequent \Delta_2, \lnot !A
    \end{align*}
    Esses são novamente exatamente os sequentes que queremos que sejam
    !!{provable} se $!C_1$ for um interpolante; tomamos de novo $!C
    \ident !C_1$. Como antes, $\Lang L(!C_1) \subseteq \Lang L(\Gamma_1,
    \Delta_1) \cap \Lang L(\Gamma_2, \Delta_2, \lnot !A)$.
    \item $\pi$ termina em $\RightR{\lif}$ com !!{formula} principal~$!A
    \lif !B$. Novamente há dois casos, conforme $!A \lif !B$ pertença à
    primeira ou à segunda parte. A !!{proof}~$\pi$ termina em
    \[
    \AxiomC{}
    \RightLabel{$\pi_1$}
    \Deduce$\maeh{!A, \Gamma_1}{\Gamma_2} \fCenter
      \maeh{\Delta_1, !B}{\Delta_2}$
    \RightLabel{\RightR{\lif}}
    \UnaryInf$\maeh{\Gamma_1}{\Gamma_2} \fCenter
      \maeh{\Delta_1, !A \lif !B}{\Delta_2}$
    \DisplayProof
    \]
    ou
    \[
    \AxiomC{}
    \RightLabel{$\pi_1$}
    \Deduce$\maeh{\Gamma_1}{!A, \Gamma_2} \fCenter
      \maeh{\Delta_1}{\Delta_2, !B}$
    \RightLabel{\RightR{\lif}}
    \UnaryInf$\maeh{\Gamma_1}{\Gamma_2} \fCenter
      \maeh{\Delta_1}{\Delta_2, !A \lif !B}$
    \DisplayProof
    \]
    Novamente atribuímos as !!{formula}s ativas, na premissa, à mesma
    parte que a !!{formula} principal na conclusão. No primeiro cenário,
    aplicamos a hipótese de indução e obtemos !!{proof}s de
    \begin{align*}
    !A, \Gamma_1 & \Sequent \Delta_1, !B, !C_1 &&\text{e} & !C_1, \Gamma_2 & \Sequent \Delta_2
    \intertext{O sequente à esquerda é uma premissa adequada para \RightR{\lif}. Assim, temos os seguintes sequentes !!{provable}:}
    \Gamma_1 & \Sequent \Delta_1, !A \lif !B, !C_1 &&\text{e} & !C_1, \Gamma_2 & \Sequent \Delta_2
    \end{align*}
    Isso significa que $!C_1$ também é um interpolante para o sequente
    final de~$\pi$, e podemos mais uma vez tomar $!C \ident !C_1$.
    
    O outro cenário é tratado do mesmo modo.
    \item $\pi$ termina em $\LeftR{\lif}$ com !!{formula} principal~$!A
    \lif !B$. Se $!A \lif !B$ pertence à primeira parte, a
    !!{proof}~$\pi$ termina em
    \[
    \AxiomC{}
    \RightLabel{$\pi_1$}
    \Deduce$\maeh{\Gamma_1}{\Gamma_2} \fCenter
      \maeh{\Delta_1, !A}{\Delta_2}$
    \AxiomC{}
    \RightLabel{$\pi_2$}
    \Deduce$\maeh{!B, \Gamma_1}{\Gamma_2} \fCenter
      \maeh{\Delta_1}{\Delta_2}$
    \BinaryInf$\maeh{!A \lif !B, \Gamma_1}{\Gamma_2} \fCenter
      \maeh{\Delta_1}{\Delta_2}$
    \DisplayProof
    \]
    A hipótese de indução nos dá agora quatro sequentes !!{provable},
    dois para o interpolante $!C_1$ da premissa esquerda e dois para o
    interpolante $!C_2$ da segunda premissa:
    \begin{align*}
    \Gamma_1 & \Sequent \Delta_1, !A, !C_1 &&\text{(1)} &
    !C_1, \Gamma_2 & \Sequent \Delta_2 && \text{(2)}\\
    !B, \Gamma_1 & \Sequent \Delta_1, !C_2 &&\text{(3)} &
    !C_2, \Gamma_2 & \Sequent \Delta_2 && \text{(4)}
    \end{align*}
    Os sequentes (1) e (3) podem servir de premissas de~\LeftR{\lif} se
    aplicarmos enfraquecimento e acrescentarmos $!C_2$ à direita de~(1)
    e $!C_1$ à direita de~(3):
    \[
    \AxiomC{}
    \Deduce$\Gamma_1 \fCenter \Delta_1, !A, !C_1, !C_2$
    \AxiomC{}
    \Deduce$!B, \Gamma_1 \fCenter \Delta_1, !C_1, !C_2$
    \RightLabel{\LeftR{\lif}}
    \BinaryInf$!A \lif !B, \Gamma_1 \fCenter \Delta_1, !C_1, !C_2$
    \DisplayProof
    \]
    Aplicando $\RightR{\lor}$, obtemos o sequente
    \[
    !A \lif !B, \Gamma_1 \fCenter \Delta_1, !C_1 \lor !C_2
    \]
    Isso sugere que $!C_1 \lor !C_2$ pode ser nosso interpolante neste
    caso. De fato, podemos !!{prove} o outro sequente exigido a partir
    dos sequentes restantes (2) e~(4):
    \[
    \AxiomC{}
    \Deduce$!C_1, \Gamma_2 \fCenter \Delta_2$
    \AxiomC{}
    \Deduce$!C_2, \Gamma_2 \fCenter \Delta_2$
    \RightLabel{\LeftR{\lor}}
    \BinaryInf$!C_1 \lor !C_2, \Gamma_2 \fCenter \Delta_2$
    \DisplayProof
    \]
    Resta verificar que $!C_1 \lor !C_2$ está na linguagem correta. A
    hipótese de indução nos dá
    \begin{align*}
    \Lang L(!C_1) & \subseteq
    \Lang L(\Gamma_1, \Delta_1, !A) \cap \Lang L(\Gamma_2, \Delta_2) &\text{e}\\
    \Lang L(!C_2) & \subseteq
    \Lang L(!B, \Gamma_1, \Delta_1) \cap \Lang L(\Gamma_2, \Delta_2)
    \intertext{Como}
    \Lang L(!A) & \subseteq \Lang L(!A \lif !B) &\text{e}\\
    \Lang L(!B) & \subseteq \Lang L(!A \lif !B)
    \intertext{temos simultaneamente}
    \Lang L(!C_1) & \subseteq
    \Lang L(\Gamma_1, \Delta_1, !A \lif !B) \cap \Lang L(\Gamma_2, \Delta_2) &\text{e}\\
    \Lang L(!C_2) & \subseteq
    \Lang L(\Gamma_1, \Delta_1, !A \lif !B) \cap \Lang L(\Gamma_2, \Delta_2)
    \intertext{e, consequentemente, $\Lang L (!C_1 \lor !C_2) = {}$}
    \Lang L(!C_1) \cup \Lang L(!C_2) & \subseteq
    \Lang L(\Gamma_1, \Delta_1, !A \lif !B) \cap \Lang L(\Gamma_2, \Delta_2),
    \end{align*}
    como exigido.

    No caso em que $!A \lif !B$ pertence à segunda parte, a situação é
    diferente, mas é tratada de modo semelhante. A hipótese de indução
    nos dá novamente quatro sequentes !!{provable}:
    \begin{align*}
    \Gamma_1 & \Sequent \Delta_1, !C_1 &&\text{(1)} &
    !C_1, \Gamma_2 & \Sequent \Delta_2, !A && \text{(2)}\\
    \Gamma_1 & \Sequent \Delta_1, !C_2 &&\text{(3)} &
    !C_2, !B, \Gamma_2 & \Sequent \Delta_2 && \text{(4)}
    \end{align*}
    Agora os sequentes (2) e (4)---com algumas !!{formula}s acrescentadas
    por enfraquecimento---servem de premissas de~\LeftR{\lif}:
    \[
    \AxiomC{}
    \Deduce$!C_1, !C_2, \Gamma_2 \fCenter \Delta_2, !A$
    \AxiomC{}
    \Deduce$!B, !C_1, !C_2, \Gamma_2 \fCenter \Delta_2$
    \RightLabel{\LeftR{\lif}}
    \BinaryInf$!A \lif !B, !C_1, !C_2, \Gamma_2 \fCenter \Delta_2$
    \DisplayProof
    \]
    Queremos novamente transformar isso em um sequente que tenha uma
    única !!{formula} construída a partir de $!C_1$ e $!C_2$, mas desta
    vez precisamos dela no antecedente (pois agora estamos tratando da
    segunda parte). Aplicar $\LeftR{\land}$ resolve o problema; devemos
    tomar $!C \ident (!C_1 \land !C_2)$ como interpolante. Felizmente,
    podemos usar os sequentes restantes (1) e (3) para !!{prove} o
    sequente correspondente para a outra parte:
    \[
    \AxiomC{}
    \Deduce$\Gamma_1 \fCenter \Delta_1, !C_1$
    \AxiomC{}
    \Deduce$\Gamma_1 \fCenter \Delta_1, !C_2$
    \RightLabel{\RightR{\land}}
    \BinaryInf$\Gamma_1 \fCenter \Delta_1, !C_1 \land !C_2$
    \DisplayProof
    \]
    Como antes, temos $\Lang L(!C_1 \land !C_2) \subseteq \Lang
    L(\Gamma_1, \Delta_1) \cap \Lang L(\Gamma_2, \Delta_2, !A \lif
    !B)$.

    \item $\pi$ termina em $\RightR{\lforall}$ com !!{formula}
    principal~$\lforall[x][!A(x)]$:
    \[
    \AxiomC{}
    \RightLabel{$\pi_1$}
    \Deduce$\maeh{\Gamma_1}{\Gamma_2} \fCenter
    \maeh{\Delta_1, !A(c)}{\Delta_2}$
    \RightLabel{\RightR{\lforall}}
    \UnaryInf$\maeh{\Gamma_1}{\Gamma_2} \fCenter
      \maeh{\Delta_1, \lforall[x][!A(x)]}{\Delta_2}$
    \DisplayProof
    \]
    ou
    \[
    \AxiomC{}
    \RightLabel{$\pi_1$}
    \Deduce$\maeh{\Gamma_1}{\Gamma_2} \fCenter
      \maeh{\Delta_1}{\Delta_2, !A(c)}$
    \RightLabel{\RightR{\lforall}}
    \UnaryInf$\maeh{\Gamma_1}{\Gamma_2} \fCenter
      \maeh{\Delta_1}{\Delta_2, \lforall[x][!A(x)]}$
    \DisplayProof
    \]
    No primeiro cenário, a hipótese de indução nos dá !!{proof}s de
    $\Gamma_1 \Sequent \Delta_1, !A(c), !C_1$ e $!C_1, \Gamma_2
    \Sequent \Delta_2$. Podemos aplicar $\RightR{\lforall}$ à primeira
    para obter $\Gamma_1 \Sequent \Delta_1, \lforall[x][!A(x)], !C_1$.
    (A condição de variável própria é satisfeita, pois já era satisfeita
    na regra \RightR{\lforall} ao final de~$\pi$.) Assim, $!C_1$ seria
    um interpolante se fosse satisfeita a condição $\Lang L(!C_1)
    \subseteq \Lang L(\Gamma_1, \Delta_1, \lforall[x][!A(x)]) \cap
    \Lang L(\Gamma_2, \Delta_2)$. Pela hipótese de indução, temos $\Lang
    L(!C_1) \subseteq \Lang L(\Gamma_1, \Delta_1, !A(c)) \cap \Lang
    L(\Gamma_2, \Delta_2)$. Precisamos, portanto, nos preocupar com~$c$:
    $c$ pode ocorrer em~$!C_1$? Não pode. Se $c$ ocorresse em $!C_1$,
    teríamos tanto $c \in \Lang L(\Gamma_1, \Delta_1, !A(c))$ (o que é
    verdade) quanto $c \in \Lang L(\Gamma_2, \Delta_2)$. Mas então $c$
    ocorreria na conclusão da inferência $\RightR{\lforall}$ em~$\pi$,
    violando a condição de variável própria.
    
    O outro cenário é semelhante.

    \item $\pi$ termina em $\RightR{\lexists}$ com !!{formula}
    principal~$\lexists[x][!A(x)]$. Novamente há dois cenários.
    Consideramos aquele em que $\lexists[x][!A(x)]$ pertence à primeira
    parte e deixamos o outro como exercício.

    No primeiro cenário, $\pi$ termina em
    \[
    \AxiomC{}
    \RightLabel{$\pi_1$}
    \Deduce$\maeh{\Gamma_1}{\Gamma_2} \fCenter
      \maeh{\Delta_1, \lexists[x][!A(x)], !A(t)}{\Delta_2}$
    \RightLabel{\RightR{\lexists}}
    \UnaryInf$\maeh{\Gamma_1}{\Gamma_2} \fCenter
      \maeh{\Delta_1, \lexists[x][!A(x)]}{\Delta_2}$
    \DisplayProof
    \]
    A hipótese de indução nos dá !!{proof}s de $\Gamma_1 \Sequent
    \Delta_1, \lexists[x][!A(x)], !A(t), !C_1$ e $!C_1, \Gamma_2
    \Sequent \Delta_2$. Podemos aplicar $\RightR{\lexists}$ à primeira
    para obter $\Gamma_1 \Sequent \Delta_1, \lexists[x][!A(x)], !C_1$.
    Assim, $!C_1$ seria novamente um interpolante se fosse satisfeita a
    condição $\Lang L(!C_1) \subseteq \Lang L(\Gamma_1, \Delta_1,
    \lexists[x][!A(x)]) \cap \Lang L(\Gamma_2, \Delta_2)$. Aqui, a
    hipótese de indução nos dá $\Lang L(!C_1) \subseteq \Lang
    L(\Gamma_1, \Delta_1, \lexists[x][!A(x)], !A(t)) \cap \Lang
    L(\Gamma_2, \Delta_2)$. Lembre-se agora de que supusemos não haver
    !!{function}s. Isso significa que o termo~$t$ só pode ser
    !!a{constant}, isto é, $t \ident b$. Se $b \notin \Lang L(!C_1)$ ou
    $b \in \Lang L(\Gamma_1, \Delta_1, \lexists[x][!A(x)]) \cap \Lang
    L(\Gamma_2, \Delta_2)$, não há motivo para preocupação: podemos
    tomar $!C_1$ como interpolante. Mas e se $b \in \Lang L(!C_1)$ e $b
    \notin \Lang L(\Gamma_1, \Delta_1, \lexists[x][!A(x)]) \cap \Lang
    L(\Gamma_2, \Delta_2)$? Em primeiro lugar, substituindo todas as
    ocorrências de~$b$ em~$!C_1$ por !!a{variable}~$y$, podemos escrever
    $!C_1 \ident \Subst{!C_1'}{b}{y}$, onde $!C_1'$ não contém~$b$.
    Além disso, ou $b \notin \Lang L(\Gamma_1, \Delta_1,
    \lexists[x][!A(x)])$, ou $b \notin \Lang L(\Gamma_2, \Delta_2)$.
    Podemos tomar, respectivamente, $!C \ident \lforall[y][!C_1'(y)]$
    ou $!C \ident \lexists[y][!C_1'(y)]$. No primeiro caso, temos:
    \[
    \AxiomC{}
    \Deduce$\Gamma_1 \fCenter
      \Delta_1, \lexists[x][!A(x)], !A(b), !C_1'(b)$
    \RightLabel{\RightR{\lexists}}
    \UnaryInf$\Gamma_1 \fCenter \Delta_1, \lexists[x][!A(x)], !C_1'(b)$
    \RightLabel{\RightR{\lforall}}
    \UnaryInf$\Gamma_1 \fCenter
      \Delta_1, \lexists[x][!A(x)], \lforall[y][!C_1'(y)]$
    \DisplayProof
    \]
    e
    \[
    \AxiomC{}
    \Deduce$!C_1'(b), \lforall[y][!C_1'(y)], \Gamma_2 \fCenter \Delta_2$
    \RightLabel{\LeftR{\lforall}}
    \UnaryInf$\lforall[y][!C_1'(y)], \Gamma_2 \fCenter \Delta_2$
    \DisplayProof
    \]
    As !!{proof}s dos sequentes superiores vêm da hipótese de indução
    (na segunda, aplicamos a admissibilidade do enfraquecimento com
    $\lforall[y][!C_1'(y)]$ no antecedente). A condição de variável
    própria de \RightR{\lforall} na primeira !!{proof} é satisfeita
    porque $b \notin \Lang L(\Gamma_1, \Delta_1,
    \lexists[x][!A(x)])$, e $!C_1'$ foi definida de modo a não
    conter~$b$.
    
    No segundo caso, temos (pela hipótese de indução, além da
    admissibilidade do enfraquecimento por~$\lexists[y][!C_1'(y)]$):
    \[
    \AxiomC{}
    \Deduce$\Gamma_1 \fCenter
      \Delta_1, \lexists[x][!A(x)], !A(b), \lexists[y][!C_1'(y)], !C_1'(b)$
    \RightLabel{\RightR{\lexists}}
    \UnaryInf$\Gamma_1 \fCenter
      \Delta_1, \lexists[x][!A(x)], \lexists[y][!C_1'(y)], !C_1'(b)$
    \RightLabel{\RightR{\lexists}}
    \UnaryInf$\Gamma_1 \fCenter
      \Delta_1, \lexists[x][!A(x)], \lexists[y][!C_1'(y)]$
    \DisplayProof
    \]
    e
    \[
    \AxiomC{}
    \Deduce$!C_1'(b), \Gamma_2 \fCenter \Delta_2$
    \RightLabel{\LeftR{\lexists}}
    \UnaryInf$\lexists[y][!C_1'(y)], \Gamma_2 \fCenter \Delta_2$
    \DisplayProof
    \]
    A condição de variável própria de \LeftR{\lexists} na segunda
    !!{proof} é satisfeita porque $b \notin \Lang L(\Gamma_2,
    \Delta_2)$, e $!C_1'$ foi definida de modo a não conter~$b$.
  \end{enumerate}
  Os outros casos são semelhantes; deixamo-los como exercícios.
\end{proof}

\begin{prob}
Suponha que tivéssemos um conectivo $\lnand$ com as regras
\[
\Axiom$\Gamma \fCenter \Delta, !A$
\Axiom$\Gamma \fCenter \Delta, !B$
\RightLabel{\LeftR{\lnand}}
\BinaryInf$!A \lnand !B, \Gamma \fCenter \Delta$
\DisplayProof
\quad
\Axiom$!A, !B, \Gamma \fCenter \Delta$
\RightLabel{\RightR{\lnand}}
\UnaryInf$\Gamma \fCenter \Delta, !A \lnand !B$
\DisplayProof
\]
Mostre que o lema de Maehara continua válido resolvendo os casos da
demonstração de \olref{lem:maehara} em que $\pi$ termina em uma dessas
regras.
\end{prob}

Uma vez estabelecido o lema de Maehara, podemos usá-lo para demonstrar o
teorema da interpolação de Craig.

\begin{proof}[Prova de \olref{thm:interpolation}]
  Suponha que $!A \Sequent !B$ seja !!{provable}. Aplique
  \olref{lem:maehara} a $\maeh{!A}{\ } \Sequent \maeh{\ }{!B}$ para
  obter o interpolante~$!C$. Temos $\Proves !A \Sequent !C$ e $\Proves
  !C \Sequent !B$.
\end{proof}

Suponha que $\Gamma$ seja uma teoria (um conjunto de !!{sentence}s) em
!!a{language}~$\Lang L$ que contenha um {predicate} $n$-ário~$R$.
Dizemos que $\Gamma$ \emph{define implicitamente}~$R$ se, e somente se,
\begin{align*}
\Gamma, \Gamma'
& \Entails \lforall[x_1][\dots\lforall[x_n][(R(x_1,\dots,x_n) \liff
R'(x_1, \dots, x_n))]],
\intertext{onde $\Gamma'$ é $\Gamma$ com todas as ocorrências de~$R$
substituídas por~$R' \notin \Lang L$. $\Gamma$ \emph{define
explicitamente}~$R$ se, e somente se, existe algum~$!A(x_1, \dots, x_n)$
que não contém~$R$ tal que}
\Gamma &\Entails
\lforall[x_1][\dots\lforall[x_n][(R(x_1,\dots,x_n) \liff !A(x_1,
\dots, x_n))]].
\end{align*}
É claro que, se $\Gamma$ define explicitamente $R$, então também define
implicitamente~$R$. A recíproca não é óbvia, mas ainda assim vale:

\begin{lem}[Teorema da definibilidade de Beth]
  Se $\Gamma$ define implicitamente~$R$, então também o define
  explicitamente.
\end{lem}

\begin{proof}
  Suponha que $\Gamma$ defina implicitamente~$R$. Sejam $R'$ e $\Gamma'$
  como acima, e seja $\Lang L' = \Lang L(\Gamma') = \Lang L \setminus
  \{R\} \cup \{R'\}$. Por hipótese, temos
  \begin{align*}
    \Gamma, \Gamma' & 
    \Sequent \lforall[x_1][\dots\lforall[x_n][(R(x_1,\dots,x_n) 
    \liff R'(x_1, \dots, x_n))]]
  \intertext{Sejam $c_1$, \dots,~$c_n$ !!{constant}s que não ocorrem em~$\Gamma$. Pelo lema da invertibilidade, temos}
    \Gamma, \Gamma', R(c_1,\dots,c_n) & \Sequent  R'(c_1, \dots, c_n)
  \intertext{Aplique o lema de Maehara a}
    \maeh{\Gamma, R(c_1,\dots,c_n)}{\Gamma'} & \Sequent
      \maeh{\ }{R'(c_1, \dots, c_n)}
  \intertext{para encontrar~$!A \ident !A(x_1, \dots, x_n)$ tal que}
    \Gamma, R(c_1,\dots,c_n) & \Sequent !A(c_1, \dots, c_n) \\
    !A(c_1, \dots, c_n), \Gamma' & \Sequent R'(c_1, \dots, c_n)
  \intertext{tenham !!{proof}s. Como $\Lang L(!A) = \Lang L_1 \cap
  \Lang L_2 = \Lang L(\Gamma) \setminus \{R\}$, $!A$ não contém~$R$.
  Substitua $R'$ por~$R$ em toda a !!{proof} do segundo sequente para
  obter !!a{proof} de}
    !A(c_1, \dots, c_n), \Gamma & \Sequent R(c_1, \dots, c_n)
  \end{align*}
  Aplique \RightR{\lif} e \RightR{\lforall} para obter !!a{proof} de
  \[
    \Gamma \Sequent \lforall[x_1][\dots\lforall[x_n][(R(x_1,\dots,x_n) \liff !A(x_1, \dots, x_n))]].
  \]
  Isso mostra que $!A(x_1, \dots, x_n)$ define explicitamente $!R$
  em~$\Gamma$.
\end{proof}

Um terceiro resultado equivalente à interpolação (e que, como ela, pode
ser demonstrado usando o lema de Maehara) é o seguinte teorema: se duas
teorias consistentes~$\Gamma_1$, $\Gamma_2$ contêm uma teoria completa
$\Gamma$ na linguagem~$\Lang L(\Gamma) \subseteq \Lang L(\Gamma_1)
\cap \Lang L(\Gamma_2)$, então $\Gamma_1 \cup \Gamma_2$ é consistente.

Lembre-se de que uma teoria completa é aquela para a qual, dada qualquer
!!{sentence}~$!A$ em sua linguagem, ou $\Gamma \Proves !A$, ou $\Gamma
\Proves \lnot !A$. Como $\Gamma_1$ e~$\Gamma_2$ são extensões
consistentes de~$\Gamma$, $\Gamma$ deve ser consistente. Segue-se que,
se $!A$ é !!a{sentence} na linguagem $\Lang L(\Gamma_1) \cap \Lang
L(\Gamma_2)$, não podemos ter $\Gamma_1 \Entails !A$ e $\Gamma_2
\Entails \lnot !A$. Para ver isso, suponha que existisse tal~$!A$. Se
$\Gamma_1 \Entails !A$, então $\Gamma \Entails !A$. Caso contrário,
pela completude de~$\Gamma$, teríamos $\Gamma \Entails \lnot !A$; como
$\Gamma \subseteq \Gamma_1$, teríamos $\Gamma_1 \Entails \lnot !A$,
tornando $\Gamma_1$ inconsistente. Portanto, $\Gamma \Entails !A$ e,
como $\Gamma \subseteq \Gamma_2$, $\Gamma_2 \Entails !A$. Logo, $\Gamma
\Entails/ \lnot !A$; caso contrário, $\Gamma_2$ seria inconsistente.

A inexistência de !!a{sentence}~$!A$ na !!{language} comum~$\Lang
L(\Gamma_1) \cap \Lang L(\Gamma_2)$ tal que $\Gamma_1 \Entails !A$ e
$\Gamma_2 \Entails \lnot !A$ é tudo o que se exige para a demonstração,
que apresentamos aqui por métodos da teoria da prova usando o lema de
Maehara.

\begin{thm}[Teorema da consistência conjunta de Robinson]
Suponha que $\Gamma_1$, $\Gamma_2$ sejam teorias consistentes e que, para
nenhum~$!A$ na !!{language} $\Lang L(\Gamma_1) \cap \Lang L(\Gamma_2)$,
tenhamos $\Gamma_1 \Entails !A$ e $\Gamma_2 \Entails \lnot !A$. Então
$\Gamma_1 \cup \Gamma_2$ é consistente.
\end{thm}

\begin{proof}
  Suponha, por absurdo, que $\Gamma_1 \cup \Gamma_2$ seja inconsistente.
  Então o sequente $\Gamma_1, \Gamma_2 \Sequent \ $ tem !!a{proof}.
  Aplique o lema de Maehara a $\maeh{\Gamma_1}{\Gamma_2} \Sequent
  \maeh{\ }{\ }$ para obter !!a{sentence}~$!A$ na !!{language}~$\Lang
  L(\Gamma_1) \cap \Lang L(\Gamma_2)$ tal que $\Proves \Gamma_1
  \Sequent !A$ e $\Proves !A, \Gamma_2 \Sequent \quad$. Desta última,
  obtemos $\Proves \Gamma_2 \Sequent \lnot !A$. Mas supusemos que não
  existe tal !!{sentence}.
\end{proof}

Acima supusemos tacitamente que as teorias $\Gamma$, $\Gamma_1$ e
$\Gamma_2$ são finitas. Pela compacidade, os resultados também valem
para teorias infinitas.

\end{document}
